\documentclass{article}
\usepackage[french]{babel}
%\usepackage[latin1]{inputenc}
\usepackage[utf8]{inputenc}  
%\usepackage[T1]{fontenc}
\usepackage[top = 3cm, left = 4cm, right = 4cm ]{geometry}
\usepackage[pdftex]{graphicx}
\usepackage{fancyhdr}
\usepackage{colortbl}
\usepackage{xcolor}

\usepackage{moreverb}

\definecolor{lightgreen}{rgb}{.5,1,.5}
\definecolor{lightred}{rgb}{1,.5,.5}

\title{
{\Huge \textbf{Rapport de projet}\\
Algorithmique et bioinformatique}}

\author{\textbf{Sneessens Joachim}\\\textbf{Tulippe-Hecq Arnaud}\\MA1 Info}


\date{Année Académique 2019-2020\\
Master 1 en sciences informatiques\\
\vspace{1cm}
Faculté des Sciences, Université de Mons}


\pagestyle{fancy}
\lhead{Sneessens J. - Tulippe-Hecq A.}
\rhead{MA1 Info}
\cfoot{\thepage}

\begin{document}

\maketitle

\vspace{9cm}

\begin{center}
\begin{tabular}[t]{c c c}
\includegraphics[height=1.5cm]{logoumons.jpg} &
\hspace{1cm} &
\includegraphics[height=1.5cm]{logofs.jpg}
\end{tabular}
\end{center}



\newpage

\section{Introduction}

%Ce projet est réalisé dans le cadre du cours d'\textit{Algorithmique et BioInformatique} donné par le Prof. Olivier Delgrange.
%donné par les enseignants Olivier Delgrange et Clément Tamines, 

L'objectif de ce projet est l'implémentation d'une solution au problème d'assemblage de fragments d'ADN. Suite au séquençage de l'ADN, un ensemble de courts fragments est obtenu, et l'assemblage consiste à aligner/fusionner ces derniers afin de retrouver la séquence d'origine. Les fragments issus du séquençage peuvent provenir du brin $3' \to 5'$, ou de son complémentaire $5' \to 3'$, ainsi lorsqu'un fragment est présent dans l'ensemble issu du séquençage, son complémentaire inversé doit aussi être considéré dans la reconstruction de la séquence originale.

Plus précisément, le but de ce projet est de fournir un programme java prenant en entrée une collection de fragments et capable de reconstituer au mieux la séquence d'origine. Au vu de la grande taille des problèmes, une grande attention doit être portée à la complexité en temps et en mémoire. Certaines collections de fragments et leurs séquences cibles respectives sont à disposition pour permettre d'évaluer la qualité du programme (la comparaison entre la séquence cible et la séquence obtenue par le 
programme étant effectuée à l'aide de l'outil $dotmatcher$).
\textcolor{red}{Précisons que l'ADN n'est vu et traité que sous son abstraction en temps que chaîne de caractères.}

Les fragments reçus en entrée sont des chaînes des caractères (dont la longueur varie entre 500 et 700 caractères). Les différents caractères sont des $a$, $c$, $t$ et $g$ correspondant \textcolor{red}{respectivement à adénine, cytosine, thymine, guanine, qui sont} au différentes molécules de base constituant l'ADN.


\section{Organisation du travail}

\textcolor{red}{
Cette section traite des différents outils utilisés pour réaliser ce projet et de la répartition des tâches au sein du groupe.}

Les différentes étapes du projet ont été implémentées dans l'ordre dans lequel elles sont présentées dans la section suivante. 
L'implémentation d'une étape ne commençait qu'une fois l'étape précédente terminée et testée. Plutôt que d'assigner à chaque membre du groupe des parties entières du travail, nous travaillions simultanément sur les mêmes parties du projet. En conséquence, il est impossible de mettre en évidence une répartition claire des tâches. Cette approche nous a permis de nous concentrer successivement sur chacune des parties du projet. S'il fallait changer quelque chose à une étape précédente, les tests unitaires nous permettaient de vérifier que les changements ne corrompent pas le programme. \textcolor{red}{ De plus, avoir tous deux connaissance de l'implémentation précise des autres parties du projet permet de m}



L'outil $git$ a été utilisé pour le contrôle de version.\textcolor{red}{, et le projet a été implémenté sous \textit{java 11}.}


\newpage

\section{\'Etapes clés}

Les différentes étapes à suivre du projet nous ont été présentées lors de la séance de présentation du projet. Le prochain paragraphe décrit brièvement leurs places respectives dans l'assemblage de fragments.

La première étape de cet algorithme consiste à construire un graphe dit \textit{overlap graph} dont les sommets sont les fragments   et les arcs sont les poids des alignements semi-globaux entre ces fragments.
Ensuite vient la construction d'un chemin hamiltonien sur ce graphe par un algorithme glouton.
L'avant dernière étape consiste à aligner tous les fragments selon l'ordre obtenu dans le chemin hamiltonien en ajoutant des gaps si nécessaire. On parle de propagation des gaps.
Finalement intervient le vote de consensus qui, d'après l'ensemble de fragments alignés obtenus après la propagation des gaps, déterminera la séquence finale que retournera le programme.

Les sections suivantes sont consacrées à une explication plus approfondie de chacune de ces étapes clés, chacune décrivant d'abord précisément le rôle de l'étape avant de s'attarder sur les choix d'implémentation propres à notre projet.

\subsection{Alignement semi-global}

L'alignement semi-global entre deux fragments est utilisé pour construire l'\textit{overlap graph}. \textcolor{red}{L'alignement semi-global vise à mesurer la ressemblance entre deux fragments sans pénaliser le décalage.} 

C'est aussi à cette étape que sont détectées les inclusions entre fragments.

Pour calculer cet alignement semi-global, on utilise la programmation dynamique.


L'arc $f \to g$ dans l'\textit{overlap graph} aura pour poids le score de l'alignement semi-global dans lequel $f$ finit par des gaps et $g$ commence par des gaps lorsqu'ils sont alignés.

Le tableau suivant montre la matrice calculée pour déterminer les alignements $f \to g$ et $g \to f$, avec $f =$ acaatgatc et $g =$ caagatcagga.

\begin{table}[h]
	\centering
	\begin{tabular}{|c|c|c|c|c|c|c|c|c|c|c|c|c|c|}
		\hline
		&& C & A & A & G & A & T & C & A & G & G & A\\		
		\hline
		& 0 & 0 & 0 & 0 & 0 & 0 & 0 & 0 & 0 & 0 & \cellcolor{lightgreen}0 & 0  \\
		\hline 
		A&\cellcolor{lightred}0 & -1 & 1 & 1 & -1 & 1 & -1 & -1 & 1 & -1 & -1 &\cellcolor{lightgreen}1 \\
		\hline 
		C&0 & \cellcolor{lightred}1 & -1 & 0 & 0 & -1 & 0 & 0 & -1 & 0 & -2 & -1 \\
		\hline 
		A&0 & -1 & \cellcolor{lightred}2 & 0 & -1 & 1 & -1 & -1 & 1 & -1 & -1 & -1  \\
		\hline 
		A&0 & -1 & 0 & \cellcolor{lightred}3 & 1 & 0 & 0 & -2 & 0 & 0 & -2 & 0 \\
		\hline 
		T&0 & -1 & -2 & \cellcolor{lightred}1 & 2 & 0 & 1 & -1 & -2 & -1 & -1 & -2  \\
		\hline 
		G&0 & -1 & -2 & -1 & \cellcolor{lightred}2 & 1 & -1 & 0 & -2 & -1 & 0 & -2  \\
		\hline 
		A&0 & -1 & 0 & -1 & 0 & \cellcolor{lightred}3 & 1 & -1 & 1 & -1 & -2 & 1  \\
		\hline 
		T&0 & -1 & -2 & -1 & -2 & 1 & \cellcolor{lightred}4 & 2 & 0 & 0 & -2 & -1  \\
		\hline 
		C&0 & 1 & -1 & -3 & -2 & -1 & 2 & \cellcolor{lightred}5 & 3 & 1 & -1 & -3  \\
		\hline 
	\end{tabular}
\end{table}

\textcolor{red}{L'alignement est reconstruit de la droite vers la gauche.}

La reconstitution de l'alignement $f \to g$ est mise en évidence en rouge dans le tableau. Pour trouver l'alignement $f \to g$, on  commence depuis la case de score maximal dans la dernière ligne. 
Partant de cette case, on va remonter le chemin vers la case $(0,0)$. Il s'agit, itérativement, de retrouver la case dont on vient parmi la case située à gauche, la case située au dessus et la case supérieure gauche.



$f \to g$
\begin{verbatim}
acaatgatc----
-caa-gatcagga
\end{verbatim}
En partant de la i\up{e} case de la dernière ligne, on sait qu'on va aligner le dernier caractère de $f$ avec le i\up{e} caractère de $g$. 

L'alignement $g \to f$ est construit en suivant les mêmes règles, mais en partant de la case de score maximal de la dernière colonne. Il est mis en évidence en vert dans le tableau.
$g \to f$
\begin{verbatim}
caagatcagga--------
----------acaatgatc
\end{verbatim}


\subsection{Overlap graph}

La première étape de cet algorithme est la construction de l'overlap graph. Il s'agit de créer un graphe dont les sommets sont les fragments donnés en entrée et leurs complémentaires-inversés. Ce graphe est orienté et pondéré. Un arc $f\to g$ de poids $w$ existe entre $f$ et $g$ si l'alignement semi-global entre $f$ et $g$ est de poids $w$.

Pour construire ce graphe, il faut donc rechercher l'alignement semi-global optimal pour chaque paire de fragments du graphe (excepté les paires constituées d'un fragment et son complémentaire-inversé).

Nous avons pris le parti de retirer du graphe les fragments inclus à d'autres.

Puisque le graphe est construit dans l'unique but de parcourir les arcs par ordre décroissant pour construire le chemin hamiltonien, nous avons choisi de mettre tous les arcs dans un tas où ils seraient triés en fonction de leur poids.

\textit{PriorityBlockingQueue}\footnote{https://docs.oracle.com/javase/7/docs/api/java/util/concurrent/PriorityBlockingQueue.html}

Pour la suite, nous nommerons $f'$ le fragment complémentaire-inversé d'un fragment $g$. 

Chaque nœud d'un graphe contient un fragment $f$ et son complémentaire-inversé $f'$. Ce choix est expliqué dans la section suivante.

Conséquemment, chaque arc du graphe est caractérisé par sa source, sa destination, son poids, un

La construction du graphe est l'étape la plus coûteuse de l'assemblage de fragments. -> multithreading grâce aux streams.



\subsection{Chemin hamiltonien}

Pour construire le chemin hamiltonien, nous allons retirer successivement sa racine au tas qui contient les arcs pour parcourir les arcs dans l'ordre décroissant. L'algorithme utilisé est très proche de l'algorithme de Kruskal, qui sert à la construction de l'arbre couvrant de poids minimal d'un graphe.
Pendant ce parcours des arcs, un arc $f \to g$ est accepté dans le chemin hamiltonien si :

\begin{itemize}
\item pour tous les arcs $i \to j$ déjà choisis dans le chemin, $i \neq f'$ et $j \neq f'$  ($f'$ étant le complémentaire-inversé de $f$)
\item pour tous les arcs $i \to j$ déjà choisis dans le chemin, $i \neq g'$ et $j \neq g'$  ($g'$ étant le complémentaire-inversé de $g$)
\item ajouter $f \to g$ au chemin ne va pas créer de cycle
\item pour tous les arcs $i \to j$ déjà choisis dans le chemin, $i \neq f$ et $j \neq g$.
\end{itemize}

L'algorithme s'arrête quand tous les sommets du graphe sont traversés par le chemin créé.

Concernant nos choix d'implémentation, nous 

\subsection{Propagation des gaps}

L'étape de propagation des gaps consiste à aligner les différents fragments selon le chemin hamiltonien construit à l'étape précédente. Cette étape consiste donc à transformer l'ensemble d'arcs du chemin en un ensemble de fragments alignés. 

Pour réaliser cet alignement de nombreux fragments, il faut, pour chaque paire d'arcs consécutifs du chemin hamiltonien (i.e., chaque paire d'arcs $f \to g$ et $g \to h$), \textcolor{red}{aligner} le $g$ obtenu dans l'alignement $f \to g$ avec le $g$ obtenu dans l'alignement $g \to h$. A chaque fois qu'un gap est inséré dans le $g$ de l'arc $f \to g$, il faut l'insérer dans tous les fragments précédant $g$ dans le chemin. Similairement, lorsqu'un gap est ajouté dans le $g$ de l'arc $g \to h$, il faut le propager dans tous les fragments suivants $g$ dans le chemin hamiltonien. 

Pour illustrer, prenons les fragments $f =$ acaatgatc, $g =$ caagatcagga et $h =$ agagtcaggacc et considérons le chemin hamiltonien $f \to g \to h$.

Voici l'alignement correspond à l'arc $f \to g$ :  
\begin{boxedverbatim}
acaatgatc----
-caa-gatcagga
\end{boxedverbatim}.

Et celui correspondant à l'alignement $g \to f$ :  
\begin{boxedverbatim}
caaga-tcagga--
--agagtcaggacc
\end{boxedverbatim}.

L'alignement $f \to g \to h$ obtenu après la propagation des gaps est celui-ci : 

\begin{boxedverbatim}
acaatga-tc------
-caa-ga-tcagga--
---a-gagtcaggacc
\end{boxedverbatim}.

Le but est donc de faire correspondre le $g$ du premier arc avec le $g$ du second en ajoutant les gaps nécessaires, tout en propageant ces gaps vers le haut ou vers le bas. On peut par exemple voir un gap ajouté à la 5\up{e} ligne et propagé vers le haut, ou encore un gap ajouté à la 8\up{e} ligne, propagé vers le bas.

En résumé, la propagation des gaps consiste à parcourir successivement les arcs du chemin, et pour chaque couple d'arcs successifs
$f \to g$ et $g \to h$, parcourir les deux $g$ obtenus de la gauche vers la droite en ajoutant des gaps en haut ou en bas dans le but d'aligner complètement ces deux $g$. Tous les gaps ajoutés sont propagés respectivement vers le haut ou le bas. Une fois que les $g$ sont complètement alignés, on passe à l'arc suivant.
\\~\\


Ici vont être abordées les spécificités propres à notre implémentation. La première question qui s'est posée au moment d'implémenter la propagation des gaps était relative à la structure de données utilisée pour stocker l'ensemble d'arcs utilisés par le chemin hamiltonien. 
L'algorithme glouton utilisé pour construire le chemin nous donne les arcs par ordre de poids décroissants, et donc pas dans l'ordre qui correspondrait à un parcours de ce chemin. Pour avoir les arcs dans cet ordre et ainsi pouvoir aligner deux à deux les arcs consécutifs, il faut trier les arcs du chemin.

\textcolor{red}{
Or pour réaliser l'alignement de tous les fragments consiste  il faut aligner tous les arcs consécutifs. Cela signifie qu'il faut, à partir d'un arc, pouvoir trouver facilement l'arc suivant (ou précédent) dans le chemin hamiltonien. Autrement}

Pour effectuer ce tri, trois solutions ont été envisagées :
\begin{itemize}
\item utiliser un arbre de recherche équilibré (\textit{TreeSet}\footnote{https://docs.oracle.com/javase/7/docs/api/java/util/TreeSet.html} de java). Cette structure garantit l'insertion et la recherche d'un élément en en $O(log(n))$. L'idée est de stocker les arcs dans l'arbre en les triant en fonction de leur source. Ainsi, une fois qu'on a trouvé le premier arc du chemin, on peut trouver les $n-1$ suivants en $O(n(log(n))$.

\item utiliser une table de hachage (\textit{HashMap}\footnote{\label{note1}https://docs.oracle.com/javase/8/docs/api/java/util/HashMap.html} de java) permettant de retrouver un arc en fonction de sa source. En utilisant une table de hachage ayant pour clés les sources des arcs et comme valeurs les arc, retrouver un arc en fonction de sa source se fait avec une complexité moyenne en $O(1)$ (bien qu'aucune garantie ne soit donnée d'avoir cette complexité pour un pire cas). Avec cette solution, trouver les $n-1$ derniers arcs du chemin se fait en moyenne en $O(n(log(n))$.

\item utiliser une table de hachage permettant de retrouver un arc en fonction de sa destination ou de sa source. Ici, l'idée est d'ajouter deux entrées dans la table pour chaque arc : une entrée associant cet arc à sa source et l'autre à sa destination. Les clés sont donc des objets constitués soit de la source et d'un booléen valant $true$, soit de la destination et d'un booléen valant $false$. Les booléens permettent de distinguer les clés qui correspondent aux sources et les clés qui correspondent aux destinations. Les valeurs sont les arcs. Cette solution permet de ne pas avoir à trouver le premier arc du chemin, puisqu'il est possible, partant d'un arc quelconque, de retrouver aussi bien l'arc suivant que l'arc précédent en $O(1)$ (en moyenne). Ainsi, il suffit de choisir un arc quelconque et trouver les $n-1$ autres arcs se fait en moyenne en $O(n(log(n))$.
\end{itemize}

C'est finalement la deuxième solution qui a été implémentée. D'après sa complexité, la première a été rapidement écartée (le pire cas des tables de hachage étant évitable grâce à un bon choix de la fonction de hachage et de la taille de la table). 
La troisième solution a comme avantage sur la deuxième de permettre de commencer à partir d'un fragment quelconque. Par contre, elle nécessite une table occupant deux fois plus d'espace et elle rend plus complexe l'implémentation. La difficulté supplémentaire d'implémentation est liée au fait que si les arcs sont traités en partant du premier, propager les gaps vers le bas revient à insérer des gaps dans un seul fragment, alors que traiter les arcs en suivant le chemin à l'envers implique de pouvoir propager les gaps vers le bas sur plusieurs fragments.

En itérant sur la liste des clés de la table qui contient le chemin, et en comparant chacune de ces clés avec l'entrée correspondante dans le tableau \textit{in} utilisé pour l'algorithme glouton, on peut trouver le premier arc du chemin en $O(n+m)$ (où $n$ est le nombre d'arcs et $m$ est la capacité de la \textit{HashMap}\footnote{\textcolor{red}{Voir la documentation java}}). En effet, la source du premier chemin est le seul nœud qui n'est pas une destination d'un arc du chemin, et donc le seul pour lequel l'entrée du tableau \textit{in} vaudra $0$. Puisque 

Au vu du peu de temps nécessaire pour trouver le premier fragment, et tenant compte des autres avantages et défauts des deuxièmes et troisièmes solutions, nous avons choisi d'implémenter la deuxième.
\\~\\ 
 
Après avoir choisi le moyen de trier nos arcs, s'est posé le problème du choix de la structure utilisée pour stocker les fragments en train d'être alignés. Pour cela, il faut une structure permettant facilement l'insertion d'un caractère à un indice donné. Les \textit{LinkedList}\footnote{https://docs.oracle.com/javase/7/docs/api/java/util/LinkedList.html} de java, qui sont des listes doublement chaînées, ont été choisies. Dans une première implémentation, naïve, une \textit{LinkedList} était utilisée pour stocker l'ensemble des caractères des fragments alignés. 
Le problème de cette implémentation est qu'elle s'est avérée extrêmement peu efficace (plus d'une demi-heure sur la plus grande collection de fragments). 

\textcolor{red}{
Ainsi, lorsqu'un gap devait être propagé vers le haut, il devait être inséré dans les $k$ \textit{LinkedList} correspondants aux $k$ fragments déjà traités. Et lorsqu'un gap devait être propagé vers le bas, il fallait l'insérer dans le fragment en train d'être traité.}

Afin d'améliorer cette première implémentation, plusieurs modifications ont été apportées.

La première d'entre elles, et la plus importante en termes d'impact sur la complexité, est la modification de la structure utilisée pour stocker les fragments en cours d'alignement (classe \textit{FragmentBuilder}). Plutôt que d'utiliser une liste chaînée pour stocker l'ensemble des caractères de chaque \textit{FragmentBuilder}, nous avons décidé d'utiliser une structure constituée de deux entiers et d'une liste chaînée. Pour illustrer les motivations ayant amené ce changement, voyons cet exemple nous montrant un alignement de 7 fragments.

\begin{verbatim}
tccgaagtctgct----------------------------------
---------tgctgctggag---------------------------
----------actactag-gcc-------------------------
----------------ag-gtcaactgatc-----------------
---------------------caactg-ccaaaaa------------
----------------------aac---ccaaaaagggg--------
-----------------------------------gggggggtcggt
\end{verbatim}

On constate que la majorité des caractères que chacun de ces fragments sont des gaps situés à la fin ou au début du mot. Plus l'exemple est grand, plus la proportion de gaps est importante. Pour le plus grand exemple, tous les mots alignés comptent environ $100000$ caractères, alors que chacun d'entre eux provient d'un fragment qui compte en $500$ et $700$ caractères. Sur les $100000$ caractères, on a donc énormément de gaps situés avant ou après une zone d'environ $700$ caractères. L'idée est d'éviter d'utiliser $n$ nœuds de \textit{LinkedList} (un Byte dans chaque nœud) pour stocker tous ces gaps.
Chaque \textit{FragmentBuilder} a donc deux entiers \textit{startGaps} et \textit{endGaps} qui stockent respectivement le nombre de gaps situés au début et à la fin du mot. Pour stocker les caractères situés entre ces gaps, une \textit{LinkedList} est toujours utilisée. Ce changement permet d'économiser de la mémoire en utilisant seulement deux entiers là où étaient nécessaires jusqu'à plus de $99000$ nœuds de \textit{LinkedList} auparavant. 
\textcolor{red}{Ce changement permet aussi d'ajouter des gaps simplement en incrémentant un entier plutôt qu'en ajoutant un noeud}

Une autre amélioration est l'utilisation de \textit{ListIterator} pour itérer sur les listes chaînées en $O(n)$. En effet atteindre le $i^e$ d'une liste chaînée demande de parcourir les $i-1$ éléments précédents (éventuellement moins dans le cas d'une liste non chaînée, mais la complexité reste inchangée). Atteindre le $i^e$ élément de la liste pour $i$ allant de 1 à $n$ donne donc une complexité du parcours en $0(n^2)$.

Le dernier changement notable est simplement l'exploitation du fait que pour deux fragments alignés $f$ et $g$, avoir $f$ situé avant $g$ dans le chemin hamiltonien (et donc $f$ situé plus haut que $g$ dans le tableau de fragments alignés) implique que \textit{endGaps}$(f) > $  \textit{endGaps}$(g)$. Ceci signifie que lorsqu'un gap est inséré à un indice se trouvant dans la zone des gaps de fin dans un fragment du tableau, alors cet indice se trouvera aussi dans la zone des gaps de fin pour tous les fragments situés plus haut dans le tableau. 


\textcolor{red}{, lors de la comparaison des deux $g_i$ des arcs $f \to g_1$ et $ g_2 \to h$, d'ajouter directement \textit{startGaps}$(g_1)$ gaps au début de $g_2$ et de comparer ensuite les deux mots caractère par caractère pour le reste du mot.}

Avec toutes ces modifications, l'étape de propagation se fait en $O(n*m)$, où $n$ est la taille du chemin hamiltonien et $m$ la taille maximale d'un fragment. En effet, on traite les $n$ fragments du chemin, et pour chacun d'eux on parcourt (en temps linéaire) uniquement la zone située entre les gaps du début et les gaps de fin. Dans la version précédente, ajouter un nouveau mot demandait d'ajouter un par un les $n$ gaps du début, et de propager sur tous les fragments déjà traités les gaps de fin. Si pour les petites instances le gain de temps n'a pas été déterminant - bien que non négligeable -, les changements ont permis de réduire le temps de propagation des gaps de plus d'une demi-heure à moins d'une demi-seconde.


\subsection{Vote de consensus}

Dernière de l'assemblage, le vote de consensus est l'étape où est construit le fragment qui sera donné en sortie. Il s'agit de retourner un fragment tel que son $i^e$ caractère est le caractère le plus présent dans la $i^e$ colonne du tableau construit pendant lors de la propagation des gaps.

Pour cette étape, notre implémentation consiste simplement à compter, pour chaque colonne, le nombre d’occurrences de chaque caractère (en ignorant les gaps) et à placer ensuite ce caractère au bon indice dans la séquence résultat.
Pour éviter d'utiliser inutilement de la mémoire en instanciant des tableaux de caractères, nous avons choisi de conserver des \textit{FragmentBuilder} pour cette étape. Pour conserver l'accès en temps constant à un élément, des \textit{Iterator} sont utilisés pour parcourir les \textit{LinkedList} des différents \textit{FragmentBuilder}.

Une petite optimisation est issue du fait que pour deux fragments alignés $f$ et $g$, avoir $f$ situé plus haut que $g$ dans le tableau \textit{FragmentBuilder} implique que \textit{startGaps}$(f) < $  \textit{startGaps}$(g)$ et \textit{endGaps}$(f) > $  \textit{endGaps}$(g)$. La première implication signifie que pour une colonne donnée, on peut s'arrêter de descendre dès qu'on arrive dans les gaps du début d'un mot, puisque chacun des mots situés plus bas aura un gap dans à cet indice. Similairement, la seconde implication permet de commencer à la ligne $j$ pour une colonne $i$ si l'indice $i-1$ se trouve dans la zone des gaps de fin de la $j^e$ ligne. 

 lorsqu'un gap est inséré à un indice se trouvant dans la zone des gaps de fin dans un fragment du tableau, alors cet indice se trouvera aussi dans la zone des gaps de fin pour tous les fragments situés plus haut dans le tableau. 

En cas d'égalité entre deux caractères, le caractère résultant est choisi selon l'ordre arbitraire suivant (ici classés par ordre de préférence décroissante) : a, t, c, g.

\section{Qualité du projet}

Nous avons eu recours aux streams pour la construction des arcs du graphe. Le multithreading est la capacité à pouvoir effectuer plusieurs taches ($threads$) simultanément. Le principal avantage est le gain de temps très important qu'apporte cette méthode. Cependant, comme les taches sont gérer en même temps , les résultats peuvent ne pas être cohérent. Le problème majeur de cette méthode est le non déterminisme des résultats.

Les tests unitaires, bien que non-exhaustifs, permettent quasiment d'avoir la certitude du bon fonctionnement de nos méthodes.



\section{Conclusion}

(+++ j'ai pas fait le blabla dans le SGA juste la table (pour l'instant je m'en occupe si tu veux)  - est ce que scinder l'intro est une bonne idée ? - est ce qu'il fut rajouter du contenu dans l'intro ? -  la répartition du travail est brève je sais aps si ça leur suffira -  j'ai refais le vote de consensus -  qualité du projet incomplet selon moi - pas d'inspi pour la conclusion +++)

utilisation de structure donné nouvelle

progression nette7

cependant - amélioration du résultat (lacune en bio)


\end{document}